\section{Unramified branching is determined by the actual vacuum}
\label{sec:sr-unramified-branching-vacuum}

The unramified multiplicity assertion needed by the two rank
recursions admits a direct proof; computing a semisimplified
Jacquet module is unnecessary here.

\begin{srlemma}
Let \(G\) be a locally compact group, \(P\) a subgroup and
\(K\) a compact subgroup such that \(G=PK\).
Let \(\pi\) be a representation generated by a vector \(v_0\)
fixed by \(K\), and let \(\tau\) be any representation of \(P\).
A \(P\)-equivariant linear map \(\pi\to\tau\) is determined
by its value on \(v_0\).
\end{srlemma}

\begin{proof}
The vectors \(\pi(g)v_0\) span \(\pi\). Write \(g=pk\);
then \(\pi(g)v_0=\pi(p)v_0\), so a \(P\)-map has value
\(\tau(p)\lambda(v_0)\) on this spanning set.
In particular, a \(P\)-map killing \(v_0\) is zero.
Applying this to the difference of two maps gives
the proportionality statement.
\end{proof}

We apply the lemma to each finite place \(v\notin T\).
Let \(G_v=U(W)(F_v)\), and let \(P_v\) stabilize the
isotropic line specified by the actual adapted plane \(H'\).
Its Levi is
\[
               E_v^\times\times U(W_1)(F_v).
\]
Let \(K_v\) be the stabilizer of the actual lattice in \(W\).
At a nonsplit place the base hypotheses make \(E_v/F_v\)
unramified with odd residue cardinality greater than three,
and the trace-dual lattice is the unimodular skew lattice
checked in the inert-model construction. Thus \(K_v\) is
hyperspecial. At a split place this is the integral
general linear subgroup after the actual lattice-basis
identification; it also covers residue characteristics
two and three.

In either case \(G_v=P_vK_v\).
Here is a lattice interpretation that also verifies the
claim for the actual line, rather than only a standard
coordinate parabolic. In the split case a line and its
compatible complementary hyperplane are represented by
primitive integral vectors and covectors with zero
pairing; elementary integral basis extension puts this
flag into standard position. In the unramified unitary
case a rational isotropic line has a primitive lattice
generator. Unimodularity supplies a partner with unit
pairing. Subtracting half its norm makes the partner
isotropic. Their hyperbolic plane is a unimodular
summand, and the integral orthogonal-complement
classification completes it to a hyperbolic lattice
basis. Consequently \(K_v\) is transitive on the
isotropic lines. This proves \(G_v=K_vP_v\);
taking inverses gives \(G_v=P_vK_v\).
All the changes here are isometries of the fixed local
ambient space.

Let \(\Pi_v,\Pi_{1,v}\) be the actual common modules at
the two ranks, as constructed in local matching. The vector
\[
                   v_0=p_v(\operatorname{vac}_v)
\]
is nonzero and \(K_v\)-fixed. At a split place this is
the normalized spherical function in the irreducible
common induced representation. At an inert place it
is the nonzero vacuum class with the common spherical
matrix coefficient. Since \(\Pi_v\) is irreducible,
\(v_0\) generates it. These properties hold for any
admissible comparison maps, not just our chosen model:
the common-action tie transports the action to their
target, and the nonzero common vacuum identifies its
irreducible spherical quotient.

For the beta step put
\[
              \chi=\beta,\qquad r=a/2;
\]
for the alpha step put
\[
              \chi=\alpha,\qquad r=a'/2.
\]
Regard \(\Pi_{1,v}\) as a \(P_v\)-module on which the
unipotent radical is trivial and the Levi acts by
\[
       (t,h_1)\longmapsto
              \chi_v(t)|t|_{E_v}^{\,r}\,\pi_{1,v}(h_1).
\]
The source predicate
\srcid{IsFiniteLocalEquivariantMap} is exactly
equivariance for this action. To check this from its
quantifiers, choose a Schwartz preimage of a vector
of \(\Pi_v\), and then choose a Schwartz preimage of
its proposed image in \(\Pi_{1,v}\). The two onto
quotient maps make both choices possible.
The predicate applied to \(u m(t,h_1)\) gives precisely
the displayed transformation rule; kernel stability
follows from the same rule applied to a zero image.
It therefore descends to an ordinary \(P_v\)-intertwiner.

If \(\lambda'\) and \(\lambda\) satisfy that predicate
and
\[
                   \lambda'(v_0)=c\,\lambda(v_0),
\]
the elementary lemma gives \(\lambda'=c\lambda\).
In particular a map vanishing on the actual vacuum
is zero. This establishes both
\srcid{UnramifiedDeterminedBySpherical} and
\srcid{UnramifiedMultiplicityOneProportional}
with the exact beta and alpha Levi characters.
It does not require that their Satake eigenvalues
be distinct, and introduces no scalar depending on
a finite character. The remaining normalization of
a ramified branching map must still be computed
on its ramified test vector.

\section{Archimedean branching from the extremal compact types}
\label{sec:sr-real-branching-compact-type}

We prove the real multiplicity assertion in the rank recursions
for the actual common Fock quotients. This supplies the source
input \srcid{ArchimedeanMultiplicityOne}, including the
correct two compact torus characters.

Work first at one infinite place. Retain the coordinates
and oscillator sign already checked in the real matching
construction, and set
\[
 \begin{gathered}
 P=a-p,\quad Q=p,\quad
 R=a+a'-(p+q),\quad S=p+q,\quad d=a+a',\\
 c=(m_\alpha-m_\beta)/2,\qquad
 u=c-(R-S)/2,\quad v=c+(R-S)/2,\\
 \lambda_0=(P-Q+m_\beta)/2,\qquad
 \mu_0=(-P+Q+m_\beta)/2 .
 \end{gathered}
\]
Here \(u,v\) are positive integers, since \(2c>d\).
The parity restrictions on the prescribed splitting make
the displayed compact highest weights integral for the
actual unitary groups. The common module for \(U(R,S)\)
contains the compact type
\[
 \tau=
 \bigl(\det_R^{\lambda_0}\otimes
                    {\mathbb S}_{(u^P)}(\mathbb C^R)\bigr)
 \boxtimes
 \bigl(\det_S^{\mu_0}\otimes
                    {\mathbb S}_{(v^Q)}(\mathbb C^S)\bigr).
\tag{1}\label{eq:sr-real-branching-type}
\]
A rectangle of height zero denotes the trivial Schur module.
This is the type constructed by the two joint-harmonic
rectangular minors. The signs and scalar shifts in (1)
are those of the actual action, not of its conjugate.

We record why its multiplicity is one in the actual
irreducible quotient. Paul--Trapa identifies that quotient,
before the common determinant twist, with the normalized
module \(A_{\mathfrak q}(\det_{P,Q}^{c}\otimes1)\), where
the ordered Levi has signatures
\((P,Q),(R-P,S-Q)\). In their normalization the inducing
weight \(\lambda\) is \(c\) on the first Levi block and
zero on the second. The smallest cross-block pairing of
\(\lambda+\rho(\mathfrak l)\) is
\[
              c-(a-1)/2-(a'-1)/2
                      =c-(d-2)/2>0.
\]
Thus it is in the good range. In the standard
cohomological-induction normalization the inducing
character is \(\lambda-\rho(\mathfrak u)\). The bottom
compact type has highest weight
\[
       \lambda-\rho(\mathfrak u)
                         +2\rho(\mathfrak u\cap\mathfrak p)
       =\lambda+\rho(\mathfrak u\cap\mathfrak p)
                            -\rho(\mathfrak u\cap\mathfrak k).
\]
Counting matrix-unit roots between the two Levi blocks
gives, after adding the common \(m_\beta/2\) determinant
shift, exactly (1). Its adjacent differences in the
positive and negative compact blocks are \(u\) and \(v\);
hence it is compact dominant. The bottom-layer theorem
therefore gives multiplicity one, since the inducing
character has multiplicity one. A precise version is
Theorem~3.2 of Pand\v zi\'c,
\emph{Dirac cohomology and the bottom layer \(K\)-types},
Glasnik Matemati\v cki 45 (2010), 453--460:
\url{https://web.math.pmf.unizg.hr/glasnik/45.2/45(2)-12.pdf}.
The theorem states that a compact-dominant bottom weight
\(\mu_L+2\rho(\mathfrak u\cap\mathfrak p)\) retains its
multiplicity from the inducing module. The normalization
conversion above verifies its hypotheses here.
The same calculation applies to the lower-rank module.

For the beta step the first signature \((P,Q)\) stays
fixed, while \((R,S)\) becomes \((R-1,S-1)\).
The lower compact type is consequently
\[
 \tau_1=
 \bigl(\det_{R-1}^{\lambda_0}\otimes
                    {\mathbb S}_{(u^P)}(\mathbb C^{R-1})\bigr)
 \boxtimes
 \bigl(\det_{S-1}^{\mu_0}\otimes
                    {\mathbb S}_{(v^Q)}(\mathbb C^{S-1})\bigr).
\]
Both rectangles fit: the beta step has \(q\geq1\),
and its adapted hyperbolic plane has both signature pieces nonzero;
hence \(S-1\geq Q\) and \(R-1\geq P\).
For the alpha step \((P,Q)\) becomes \((P-1,Q-1)\)
and \((R,S)\) again becomes \((R-1,S-1)\);
the difference \(P-Q\) and the integers \(u,v\) stay
unchanged. Its lower type is
\[
 \tau_1=
 \bigl(\det_{R-1}^{\lambda_0}\otimes
                    {\mathbb S}_{(u^{P-1})}
                                             (\mathbb C^{R-1})\bigr)
 \boxtimes
 \bigl(\det_{S-1}^{\mu_0}\otimes
                    {\mathbb S}_{(v^{Q-1})}
                                             (\mathbb C^{S-1})\bigr).
\]
Here \(p\geq1\) and \(q=0\), so \(P,Q\geq1\) and all
the indicated rectangles again fit. Each \(\tau_1\)
has multiplicity one in its actual lower quotient.

Choose the Cartan-compatible frames of the actual adapted
plane \(H'\). A compact isometry sends its isotropic line
to \(\mathbb C(e_R+f_S)\), with the orthogonal lower space
in the remaining \(R-1,S-1\) coordinates. This follows
by decomposing an isotropic vector into its positive
and negative parts, which have equal norm, and using
the two compact unitary groups to put those parts into
their last coordinates. Adaptedness makes the lower
space Cartan stable, so these are also its prescribed
compact coordinates. The compact part of the line
parabolic is
\[
 U(1)\times U(R-1)\times U(S-1),
\]
where \(z\in U(1)\) multiplies both removed coordinates
\(e_R\) and \(f_S\).

We need only the elementary rectangular branching rule.
On restriction
\[
 {\mathbb S}_{(u^P)}(\mathbb C^{R-1}\oplus\mathbb C)
          \quad\hbox{to }U(R-1)\times U(1),
\]
each summand has multiplicity one, and the exponent
\(e\) on the last \(U(1)\) coordinate lies in
\(0\leq e\leq u\). The \(e=0\) summand, when \(P<R\),
is \({\mathbb S}_{(u^P)}(\mathbb C^{R-1})\);
the \(e=u\) summand, when \(P\geq1\), is
\({\mathbb S}_{(u^{P-1})}(\mathbb C^{R-1})\).
One proof is the semistandard-tableau branching rule:
the boxes bearing the last label form a horizontal
strip. A rectangle permits at most one box in each
of its \(u\) columns; removing none or all of that
strip gives exactly the two stated extreme shapes.
Every intermediate interlacing shape determines its
strip and occurs once. If \(P=0\) only \(e=0\)
occurs, and if \(P=R\) only \(e=u\) occurs.
Use the identical statement with \(v,Q,S\), writing
\(f\) for its last-coordinate exponent.

On the compact part of the line Levi, (1) therefore
has torus exponent
\[
                 \lambda_0+\mu_0+e+f
                          =m_\beta+e+f .
\]
The exact beta character is \(\beta_\infty(z)=z^{m_\beta}\);
its absolute-value factor is one on \(U(1)\).
Hence it forces \(e=f=0\). Its permitted lower compact
type is exactly the beta \(\tau_1\) and occurs once.
The exact alpha character is
\(\alpha_\infty(z)=z^{m_\alpha}
                  =z^{m_\beta+u+v}\).
It forces \(e=u,f=v\), again giving exactly its
\(\tau_1\), once. This argument also covers the boundary
\(Q=S\) in the alpha step: then \(f=v\) is already
forced by rectangular branching. In particular no
unproved uniqueness statement about Jacquet parameters
or lowest-type labels enters this restriction calculation.

Finally the source's real local-equivariance predicate
says precisely that the linear map intertwines the Lie
algebra of this real parabolic, with its specified Levi
character, and the compact parabolic subgroup. The
unipotent condition kills its nilpotent Lie algebra;
the Levi derivative condition gives the infinitesimal
character rule; the compact condition gives the torus
character just computed.

An irreducible Harish--Chandra module is generated by
any nonzero compact type, so \(\Pi=U(\mathfrak g)\tau\).
The real Iwasawa decomposition gives
\(\mathfrak g=\mathfrak p_{\rm line}+\mathfrak k\)
after complexification. Ordering a basis and applying
the Poincar\'e--Birkhoff--Witt theorem yields
\[
                 \Pi=U(\mathfrak p_{\rm line})\tau .
\]
Thus a parabolic-equivariant map is determined by its
restriction to \(\tau\). The compact branching calculation
bounds that restriction space by the multiplicity of
\(\tau_1\) in \(\Pi_1\), namely one. The restriction map
is injective by the displayed spanning formula.
It follows that the space of real branching maps has
dimension at most one. If one is nonzero, every other
one is its scalar multiple, exactly as required.

For all infinite places take the tensor product of these
compact types. The compact torus at each place is
independent, so the two extreme-exponent conditions
hold place by place. Their restriction space still
has dimension one, and the product parabolic
Poincar\'e--Birkhoff--Witt argument is unchanged.
This establishes the source input for the full
archimedean tensor module and any actual comparison
quotients. It identifies the branching line, without
normalizing an arbitrary vector on that line. In an
actual rank recursion the real evaluation and cone
operators must still be shown to be nonzero and
normalized using only the fixed infinite data.
That last normalization cannot follow from
one-dimensionality alone.
